\paragraph{Keeping the selected mode under a change of frame.}

To construct the CM arc, follow one nonzero Fourier coefficient through the change of flag.

Fix the component and a nonzero low coefficient \(F_\nu\) on a
rational isotropic line \(E\xi\). Let \(a_\tau(x)>0\) be the
cusp time at a flag frame \(x\), normalized as the reciprocal of
\(c_\tau\operatorname{maj}_{x,\tau}(\xi)\), with \(c_\tau\) the
positive flow constant of the component. The complex transvection
with parameter \(s=(s_\tau)_\tau\) sends these times to
\(a_\tau(x)+\operatorname{Im}s_\tau\). Its domain is the product
of half-planes given by positivity of the translated times.
Real transvections act on the Fourier coefficient by

\[
  \exp(4\pi i\sum_\tau\nu_\tau s_\tau).
\]
The product of the flowed coefficient and the reciprocal
exponential is holomorphic on these half-planes, by integrating
the section over its compact period torus. It is constant under
real translations. Its complex derivatives in those same
directions vanish, so the product is constant throughout the
domain. In particular transporting a nonzero reading to all
times equal to one keeps it nonzero.

For a flag write \(\widehat N(x)\) for the supremum of the
normalized point norms of \(F_\nu\) over its frames. This is
finite by the Fourier coefficient bound on the compact real
period torus, since those real transvections preserve the point
metric. It is positive at some flag of time one. Put

\[
  S=\sup_{a(x)=1}\widehat N(x)>0.
\]
 Take a flag of time one
whose supremum is greater than \(S/2\), then a frame on it
whose point norm is still greater than \(S/2\).
Signed Gram--Schmidt is a right upper triangular change of frame.
The weight character of the section and the norm factor cancel
under this change, and the majorant depends only on the signed
flag. We may therefore take an orthonormal frame \(x_0\) of the
same flag, with point norm between \(S/2\) and \(S\).

\begin{proposition}[A split frame detecting the selected mode]
\label{prop:split-frame-low-mode}
Fix the torsion order before the definite lift. If \(a_\nu\) is a nonzero selected low coefficient, it is detected by a bounded split approximation,
\[
 a_\nu\ne0\quad\Longrightarrow\quad \exists\,x_{\rm split}\quad a_\nu(x_{\rm split})\ne0.
\]

The construction keeps the component and isotropic line fixed.
It chooses the quotient complement and the frame after the coefficient
and order are fixed, using the metric and radius bounds below.
\end{proposition}

\begin{proof}
\proofstep{A motion that preserves a large coefficient.}
Put \(\rho=(12\pi(c_1+w_0))^{-1}\).
For a frame with times in
\([(1-\rho)/(1+\rho),(1+\rho)/(1-\rho)]\), move its times to one
as above. A bound for the ratio of the norm factors follows
from the leading Gram minors under that imaginary transvection.
The change in the Gram matrix under a complex transvection
with imaginary time \(s_\tau\) is a Hermitian rank-one
update. Its coefficient has the sign of \(s_\tau\).

For the \(i\)-th leading block let \(G_i(0)\) and \(G_i(s_\tau)\) be its Gram matrices before and after the update. The determinant lemma gives
\[
  \frac{\det G_i(s_\tau)}{\det G_i(0)}=1+s_\tau b_i.
\]
 Here \(b_i\) is the signed squared projection
of the line onto that block in the flow normalization.
Since the line is isotropic, the leading coefficients start
and end at zero, increase on the positive block and decrease
on the negative block, with maximum \(1/a_\tau\).
Thus the ratios lie between \(1\) and \(1+s_\tau/a_\tau\)
in either order.
More explicitly, with \(a_\tau=a_\tau(x)\), the norm factor at
\(x\) divided by the norm factor after transport is bounded by

\[
  \prod_\tau\max(a_\tau,a_\tau^{-1})^{s_\tau/2},
\]
 where
\(s_\tau\) is the spread of the weight at \(\tau\). Both time
ratios are at most \(1+3\rho\). This contributes at most
\(\exp(3\rho w_0/2)\le e\).
The quotient of the raw coefficient norms is
\(\exp(4\pi\sum_\tau \nu_\tau(1-a_\tau))\).
Since \(\nu_\tau\ge0\), \(\sum_\tau\nu_\tau\le c_1\) and
\(1-a_\tau\le2\rho\), the quotient is at most

\[
  \exp(8\pi\rho c_1)\le e.
\]
 Thus
\begin{equation}\label{eq:slice-split-near-sup}
  \widehat N(x)\le e^2 S .
\end{equation}
The movement and the comparison apply to every frame of the flag,
which explains the supremum on the left.

\paragraph{The first minimum on the isotropic quotient.}
Use \(L'\) for the Hermitian self-dual lattice of the component,
and \(\mathcal L\) for its period lattice. The fixed factors satisfy
\(N_A\mathcal L\subset L'\) and \(N_C L'\subset\mathcal L\). The
majorant-orthogonal complement of the hyperbolic plane at \(x_0\)
represents \(\xi^\perp/E\xi\). Write \(v_\tau\) for the image there
of a putative short nonzero class represented by \(v\in L'\cap\xi^\perp\).
Let \(N_0,H_0,K_0',R_0\) be the split parameters defined below, and
\(e_{\rm sp}=\min(1,f^2/(N_0^2 N_C^2 c_1 R_0))\).

We claim that
\[
 \max_\tau\operatorname{maj}_{x_0,\tau}(v_\tau)\ge e_{\rm sp}.
\]
If the maximum is smaller, choose \(0<\epsilon_0<e_{\rm sp}\)
bounding it. The Hermitian self-pairing of \(v\) is integral
and is the pairing of \(v_\tau\) at every embedding. It has
absolute value less than one everywhere, so it is zero.
Shift \(v\) by \(N_A b\xi\), \(b\in\mathcal O_E\), to make its
majorant projection coefficient \(b_\tau'\) onto \(\xi\) satisfy

\[
  |b_\tau'|\le 2f^2\operatorname{rd}(E)N_A
\]
 at all places.
Indeed the Euclidean dual of the embedded \(\mathcal O_E\)
is the conjugate codifferent scaled by two, by the trace formula.

For each nonzero element of the codifferent the field norm has
absolute value at least \(1/d_E\), by the inclusion of its
principal ideal in the codifferent. The arithmetic--geometric
mean inequality gives Euclidean dual minimum at least
\(2\sqrt f/\operatorname{rd}(E)\).
Let \(\rho_2\) be the Euclidean covering radius and \(\lambda_1^*\) the dual minimum. Banaszczyk's bound in dimension \(2f\) gives
\[
  \rho_2\le\frac{2f}{2\lambda_1^*}\le\frac{\sqrt f\,\operatorname{rd}(E)}2.
\]
The maximum norm radius is therefore at most
\(2f^2\operatorname{rd}(E)\). \cite{banaszczyk}
The shift leaves \(v_\tau\) and the assumption of shortness
unchanged. The two vectors \(\xi,N_Cv\) are now in \(\mathcal L\)
and span a totally isotropic plane.

\proofstep{The average along the transvection torus.}
Average along the real transvection torus for this plane. Its
characters are Hermitian two-by-two frequency matrices
\(\chi_\tau=\tau(\beta)\). The principal level torus contains the
translations obtained from integral Hermitian matrices after
multiplying by the level. Trace duality, including the off-diagonal
trace pairing, gives 
\[
  N_0\beta_{ij}\in\mathcal O_E,
\]
 with
\(N_0=(g_0g_1)d_E^2\). Taking the \(\xi\)-coefficient first retains
just \((\chi_\tau)_{00}=\nu_\tau\). The surviving frequencies are
positive definite at every embedding.

In fact translating in a
rank-one positive direction with negative frequency would give
exponential growth against the polynomially varying norm factor.
This gives positive semidefiniteness. If one determinant vanishes,
the algebraic determinant vanishes. A rational vector in the
kernel gives a rational isotropic line on which the frequency
is zero. The corresponding constant-term projection is zero by
cuspidality. Translation averages commute on this abelian torus,
so such a frequency has zero coefficient.

Normalize the direction \(v_\tau\) to have cusp time one, writing
it as \(a_{\tau,0}\xi+a_{\tau,1}N_Cv\) after this rescaling.
The normalization gives
\[
 |a_{\tau,1}|^2
   = (N_C^2 c_\tau\operatorname{maj}_{x_0,\tau}(v_\tau))^{-1},
 \qquad 0<c_\tau\le2 .
\]
Write \(d_\tau=\det\chi_\tau>0\), let \(R_\tau\) be the value of
this frequency on the normalized direction \(a_\tau\), and put
\(T=2\sum_\tau R_\tau\). Completing the square in the first
coordinate gives \(\nu_\tau R_\tau\ge d_\tau|a_{\tau,1}|^2\).

The algebraic integer \(\det(N_0\beta)\) is nonzero, so the
product of \(N_0^2d_\tau\) over the complex half is at least one.
Since \(\sum_\tau\nu_\tau\le c_1\), two applications of the
arithmetic--geometric mean inequality give
\[
 \sum_\tau d_\tau/\nu_\tau\ge \frac{f^2}{N_0^2c_1},
 \qquad T\ge\frac{f^2}{N_0^2N_C^2c_1\epsilon_0}> R_0 .
\]
The first inequality uses
\(\prod_\tau(f\nu_\tau/c_1)\le1\). In particular all the
\(\nu_\tau\) under consideration here are positive.

\proofstep{The first minimum at the same frequency.}
The height count is on the same frequencies. In each embedding
write the second standard vector as \(q a_\tau+p(1,0)\).
Then \[
 |q|^2=N_C^2 c_\tau\operatorname{maj}(v_\tau)\le2N_C^2,\qquad |p|^2=N_C^2|b_\tau'|^2.
\]
Positivity and Cauchy--Schwarz bound the second diagonal and
the off-diagonal frequency entries by \(H_0(c_1+T/2)\).
For the off-diagonal use

\[
  |(\chi_\tau)_{01}|^2\le\nu_\tau(\chi_\tau)_{11}.
\]
The first diagonal is prescribed, so the other two entries
determine the frequency. In the maximum of the complex norms
the integer lattice has minimum at least one, by the field norm.
Disjoint translates of the open ball of radius \(1/2\) give at
most \((1+2B)^{2f}\) lattice points in a ball of radius \(B\).
For \(t\ge1\) use \(B=N_0 H_0(c_1+t/2)\) for each of the two
entries. Counting both entries therefore bounds the number of Hermitian frequency matrices \(\beta\) by
\[
  \#\{\beta:T(\beta)\le t\}\le(K'_0t)^{4f}\qquad(t\ge1).
\]

The normalized \(v_\tau\)-transvection commutes with the
\(\xi\)-flow. The time-one flag, and every frame in it, can be
moved back by \(i/2\) in the normalized \(v_\tau\) direction
at each place. Its time becomes \(1/2\), and the \(\xi\)-time
stays one, since the two directions are majorant-orthogonal.
The real period averages of the rank-two group at that moved
point therefore bound the normalized coefficient by \(S\).
Transport back multiplies the raw coefficient by \(e^{-\pi T}\)
and changes the norm factor by at most \(2^{w_0/2}\).
Group the resulting bounds by \(m=\lceil T\rceil\). For
\(m>R_0=w_0+4f\log(8fK_0')+4f+2\), the inequality
\(\log(m/(8f))\le m/(8f)-1\) gives

\[
  (K_0'm)^{4f}\le e^{2m-w_0-2}.
\]
Thus the sum of the coefficient bounds, divided by \(S\), is
at most
\[
 2^{w_0/2} e^{\pi-w_0-2}
 \frac{e^{4(2-\pi)}}{1-e^{2-\pi}}<\tfrac12 .
\]
The absolute bound justifies recovering the section from its
Fourier series on the compact torus. Taking the supremum over
the original flag contradicts the choice above \(S/2\). This
proves the quotient minimum with the cutoff fixed at \(c_1\).
\label{cor:slice-split-frame-1}

\paragraph{A signed rational complement.}

The quotient on the isotropic line controls a first minimum. To turn that estimate into a CM frame we need a rational complement with the prescribed signs.

The quotient of \(L'\cap\xi^\perp\) by \(L'\cap E\xi\) is
a self-dual Hermitian lattice on \(\xi^\perp/E\xi\).
One can check this locally. A generator of the saturated line
is primitive, and self-duality supplies a partner pairing to one.
If these vectors are \(\xi_0,\eta_0\), the Gram
matrix on their plane is
\[
 \begin{pmatrix}0&1\\1&H(\eta_0,\eta_0)\end{pmatrix}
\]
and its determinant is \(-1\). Its inverse is integral,
so subtracting the plane projection preserves the lattice.
Inside this plane the perpendicular to \(\xi_0\) is
the line it spans. Thus the quotient lattice is the
integral orthogonal complement. Its signature in the selected convention
is \((q_j-1,r-q_j-1)\).

\proofstep{Choosing the rational complement by transference.}
Use the transference bound for this lattice with lower majorant
\(e_{\rm sp}\). In real dimension \(2f(r-2)\) it gives covering
radius and an independent short family at

\[
  \mu=4(f(r-2))^3d_E/\sqrt{e_{\rm sp}}.
\]
Use the same lattice construction as for the full interior frame
with approximation allowance \(\rho/40\). Namely, approximate
the positive axes scaled by \(4\sqrt{r-2}\mu/(\rho/40)\),
and multiply the orthogonal projections of the short spanning
family to the negative space by the positive Gram determinant.
The positive block errors have operator norm at most
\(\sqrt{r-2}\mu\), so its span is a graph of norm at most
\(\rho/120\). The negative projection is given by the Gram
adjugate and is integral.

The leading positive Gram determinant
and the Gram-adjugate formula give the length bound
\(B_{\rm rad}^{2q_j-1}\) on the resulting letters, with
\(B_{\rm rad}=2\cdot4\sqrt{r-2}\mu/(\rho/40)\).
Indeed, for a vector \((u,v)\) the negative projected norm
squared is
\[
 \|(1-AA^*)^{-1/2}(v-Au)\|^2\le (1+\|A\|^2)\|(u,v)\|^2/(1-\|A\|^2)\le2\|(u,v)\|^2.
\]
Raising the common bound
for the positive lengths to twice the number of positive
axes pays the determinant. If \(q_j=1\) there is no positive
axis in the quotient, and the short negative family needs
no projection. Lift the letters to \(L'\cap\xi^\perp\).
Their pairings and their quotient basis are independent
of the choices of lifts.

In the signed orthonormal basis put
\(C_+=(1-A^*A)^{-1/2}\) and \(C_-=(1-AA^*)^{-1/2}\).
The block matrix with the required positive graph is
\[
                 U_A=\begin{pmatrix}C_+&A^*C_-\\
                                        AC_+&C_-\end{pmatrix}.
\]
For the form with matrix \(\operatorname{diag}(1,-1)\),
the two column blocks are orthogonal: their mixed matrix
is \(C_+(A^*-A^*)C_-=0\). Their squared Gram blocks are
\[
 C_+(1-A^*A)C_+=1,\qquad C_-(AA^*-1)C_-=-1.
\]
Thus \(U_A\) is unitary. Extend it by the identity on the
hyperbolic plane.

The
resulting frame \(x_1\) has the selected letters in the two
eigenspaces on the complement, and its \(\xi\)-time is still
one. The positive plane in the chart of \(x_0\) has graph
\(Z_1\) of norm at most \(\rho/100\). With
\(t=(1-\rho/100)/(1+\rho/100)\), the majorant at \(x_1\) is
between \(t\) and \(1/t\) times that at \(x_0\). The quotient
lower bound for \(x_1\) is therefore \(t e_{\rm sp}\).

\paragraph{Clearing the projection at the prescribed order.}
Let \(W_\tau\) denote the complex complement at \(x_1\),
and write \(P_\tau(\mu)\) for the coefficient of \(\xi\)
in the projection of \(\mu\in\mathcal L\cap\xi^\perp\).
The coefficient projection comes from the exact sequence
\[
 0\longrightarrow \mathcal L\cap E\xi
 \longrightarrow \mathcal L\cap \xi^\perp
 \longrightarrow
   (\mathcal L\cap\xi^\perp)/(\mathcal L\cap E\xi)
 \longrightarrow 0.
\]
The quotient is torsion-free because the submodule is the
intersection with an \(E\)-line. It is finitely generated over
the Dedekind ring \(\mathcal O_E\), hence projective, so the
sequence splits over \(\mathcal O_E\). Write \(l(\mu)\xi\) for an integral
projection from such a splitting. Its coefficient extends
by scalar extension to an \(E\)-linear form with 
\[
  l(\xi)=1.
\]
Hence the difference \(P_\tau-l_\tau\) is a linear form on
the quotient.

\proofstep{Simultaneous clearing at the prescribed torsion order.}
We approximate \(n'(P_\tau-l_\tau)\) simultaneously by the
linear forms of one lattice element. In the real tangent
space \(\prod_\tau W_\tau\) with the Euclidean majorant metric,
the image of \(\mathcal L\cap\xi^\perp\) is a lattice of
dimension \(d=2f(r-2)\). The quotient minimum is at least
\(\sqrt{t e_{\rm sp}}/N_A\): multiplication by \(N_A\) sends
the projected lattice to the one used for the first-minimum bound.
The marked Riemann pairing is
\(\operatorname{Tr}_{E/\mathbb Q}(-\xi_0 H(-,-))\) on period
vectors, where 
\[
  2|\tau(\xi_0)|=k_0.
\]
This pairing is perfect on the projected lattice, since it is
perfect on the period lattice and the isotropic sublattice is
saturated.

To see this, extend a basis of the primitive
isotropic sublattice by symplectic partners, obtained by lifting
their dual functionals under the perfect pairing. The resulting
even-rank block is unimodular, and the perpendicular quotient
is its integral orthogonal complement.
Relative to the majorant, the real operator associated to the
pairing is \(k_0\) times an orthogonal map.
Writing \(T=k_0 U\) with \(U\) orthogonal for the majorant gives
\[
  \|Tx\|=k_0\|x\|,\qquad T^*T=k_0^2I .
\]
 Indeed it is the
product of multiplication by a purely imaginary scalar,
the signature involution, and the factor two from the CM trace.

The similarity \(T\) takes the projected lattice to its Euclidean dual.
Thus the dual minimum is at least \(k_0\sqrt{t e_{\rm sp}}/N_A\).
Transference and rounding in \(d\) independent vectors give
covering radius at most
\(d^3N_A/(2k_0\sqrt{t e_{\rm sp}})\) for the projected lattice.
Via Riesz representation, the linear forms
\[
 u\longmapsto d_E(-\tau(\xi_0)) H_\tau(u,\mu_{0,\tau}),
 \qquad \mu_0\in\mathcal L\cap\xi^\perp ,
\]
come from the projected lattice by an operator of norm
\(d_Ek_0/2\) times an isometry. The covering estimate
therefore approximates the tuple \(n'(P_\tau-l_\tau)\)
with error in its representing vector bounded by
\(d_Ed^3N_A/(4\sqrt{t e_{\rm sp}})\).
The majorant of \(\xi\) is \(1/k_0\).

Multiplying the error
norm by \(1/\sqrt{k_0}\) for its value on the line is
covered by \(B_{\rm ind}(t e_{\rm sp})\), the index allowance
below evaluated with lower bound \(t e_{\rm sp}\).
For the discriminant factor, the pairings \(H(L',L')\)
generate the unit ideal by Hermitian self-duality.
Since 
\[
  N_C L'\subset\mathcal L,
\]
 integrality of the trace
pairing puts \(N_C^2(-\xi_0)\) in the codifferent. All its
embedding norms equal \(N_C^2 k_0/2\). The ideal norm lower
bound \(1/d_E\) therefore gives
\(N_C^2 k_0\operatorname{rd}(E)\ge2\), as used in the index
allowance.
All these comparisons concern the projected lattice of the
full period module, independent of the particular spanning
letters just chosen.

\proofstep{Returning to the Hermitian frame.}
Dividing by \(n'\) gives a correction \(\delta_\tau\) on \(W_\tau\)
of operator norm, as a map into the line with its majorant metric,
at most \(\kappa=B_{\rm ind}(t e_{\rm sp})/n'\).
It has the required integral effect. For \(\mu\in\mathcal L\cap\xi^\perp\)
the corrected projection multiplied by \(n'\) has coefficient
\[
 n'l(\mu)+d_E(-\xi_0)H(\mu,\mu_0).
\]
The second summand is integral since its value before multiplication
by \(d_E\), paired against each \(a\in\mathcal O_E\) by the trace,
is \(\operatorname{Tr}_{E/\mathbb Q}(-\xi_0 H(a\mu,\mu_0))\in\mathbb Z\).
Thus it is in the trace dual. The first summand belongs to the
line coefficient module by the integral splitting. The vector
with the displayed coefficient is in \(\mathcal L\cap E\xi\).

At each embedding choose \(\xi'\) in the hyperbolic plane with
\(H_\tau(\xi,\xi')=1\), of norm squared the reciprocal of
\(\operatorname{maj}_{x_1,\tau}(\xi)\), and isotropic.
Take \(w_\tau\in W_\tau\) with \(H_\tau(u,w_\tau)=\delta_\tau(u)\).
The unitary fixing \(\xi\) acts by
\[
 u\longmapsto u+\delta_\tau(u)\xi,\qquad
 \xi'\longmapsto
 \xi'-w_\tau-\tfrac12 H_\tau(w_\tau,w_\tau)\xi .
\]
Preservation of all pairings follows by expanding them on the
three summands. In this moved frame the line projection is
\(P_\tau-\delta_\tau\) and the eigenspaces on the quotient are
unchanged. The two linear off-diagonal terms have operator
norm at most \(\kappa\) each and the quadratic term at most
\(\kappa^2/2\), so the move differs from the identity by at
most \(3\kappa\) in the \(x_1\) majorant when \(\kappa\le1\).

In the \(x_0\) majorant the bound is \(\eta=3\kappa/t\).
With \(z_1=\rho/100\), the transformed graph has norm at most
\[
 \frac{z_1+\eta(1+z_1)}{1-\eta(1+z_1)}<\rho/21 .
\]
Indeed the new upper block on \(\binom{1}{Z_1}\) differs from
the identity by at most \(\eta(1+z_1)\), and the lower block
changes by at most the same amount. The inequality on \(n'\)
gives 
\[
  \kappa\le\rho/(1000t),
\]
 and \(\rho\le1/10\) gives the
strict numerical bound. This construction uses the \emph{chosen}
positive integer \(n'\). Clearing denominators first and then
taking an arbitrary multiple would give a different quantifier.
\label{cor:slice-split-frame-2}

\paragraph{The small-chart nonvanishing test.}
Let \(Z\) be the lower-block chart parameter of a final moved
flag in the orthonormal frame \(x_0\). Use the operator norm on
each block, and the maximum over the complex places. For
\(\|W\|\le\rho\) the points \(x(W)\) of this chart are in the
domain. For each singular value \(s\) of \(W\), the majorant in the
associated signed two-plane has eigenvalues \((1+s)/(1-s)\)
and its reciprocal relative to \(x_0\).
The two relative eigenvalues are reciprocal, as can be read from
\[
  \frac{(1+s)^2}{1-s^2}=\frac{1+s}{1-s},\qquad \frac{(1-s)^2}{1-s^2}=\frac{1-s}{1+s}.
\]
 This follows from the
unitary graph matrix above, whose two-by-two block is
\((1-s^2)^{-1/2}\left(\begin{smallmatrix}1&s\\s&1\end{smallmatrix}\right)\).
The remaining eigenvalues are one. Thus the majorants change by
factors between \((1-\rho)/(1+\rho)\) and its reciprocal.

For a dominant weight
of spread at most \(w_0\), the norm factor \(J(x)\) satisfies
\[
  |\log J(x(W))-\log J(x_0)|\le 4w_0 r\|W\|^2\le \tfrac14 .
\]
For leading minors ending in the positive block one takes
determinants of \(1-W_1^*W_1\), with \(W_1\) the relevant columns
of \(W\). For larger minors, eliminate the retained negative
identity block. The remaining factor is \(1-W_2^*W_2\), with
\(W_2\) the rows of \(W\) after deleting that block. The full
determinant has absolute value one.

For each of the other positive absolute determinant ratios \(q_i\), the block bound gives
\[
  |\log q_i|\le-r\log(1-\|W\|^2).
\]
Summation by parts in the consecutive minor ratios charges
only the weight spread, giving the stated quadratic estimate.
This and the majorant comparison put the chart in the range of
\eqref{eq:slice-split-near-sup}. If \(\|Z\|<\rho/21\), put
\(R=\rho/\max(\|Z\|,\rho/22)>21\).
The function \(g(t)=J(x_0)F_\nu(x(tZ))\) on \(|t|\le R\)
has \(|g(t)|\le e^{9/4}S\), and \(|g(0)|>S/2\).
If \(g(1)=0\), precompose with the automorphism of this disc
that exchanges 0 and 1. Schwarz's lemma gives

\[
  |g(0)|\le e^{9/4}S/R<S/2,
\]
 a contradiction.
This is the form of the nonvanishing argument used in the split
detector. It allows the frame move to depend on the coefficient
and the prescribed torsion order.
\label{cor:slice-split-frame-3}
\end{proof}
